
\chapter{Summary in French}

\startcontents[chapters]
\color{blue}

\color{black}
L'attention est un domaine amplement discut\'e et \'etudi\'e en neurosciences, en psychologie, en sciences cognitives et en apprentissage automatique \citep{chun2011taxonomy,cho2015describing}. L'attention est le processus consistant \`a se concentrer de manière s\'elective sur un aspect discret de l'information tout en ignorant les autres informations perceptibles. Une caract\'eristique largement accept\'ee de l'attention est qu'elle facilite l'utilisation efficace des ressources informatiques disponibles.
Bien que l'attention ait \'et\'e \'etudi\'ee depuis des d\'ecennies, elle est encore loin d'\^etre un concept simple ou unifi\'e \citep{lindsay2020attention}.

La litt\'erature en sciences cognitives d\'epeint plusieurs aspects de l'attention, comme le fait qu'elle puisse \^etre concentr\'ee, focalisée sur une modalit\'e particuli\`ere, divis\'ee, \^etre s\'elective et avoir une capacit\'e finie. Cependant, la s\'electivit\'e reste son trait le plus caract\'eristique. La s\'electivit\'e est n\'ecessaire en raison de la disponibilit\'e limit\'ee des ressources. R\'ecemment, l'attention visuelle a fait l'objet d'un intérêt consid\'erable dans le domaine de l'intelligence artificielle. L'attention visuelle \citep{ahmad1991visit} est la capacit\'e \`a hi\'erarchiser les informations tout en n\'egligeant les informations non pertinentes pour contenir la surcharge de donn\'ees dans notre syst\`eme visuel. L'attention visuelle permet de r\'epondre \`a la question: \textit{quoi} regarder et \textit{où} regarder. Cela a \'et\'e largement \'etudi\'ee en psychologie et en neurosciences \citep{posner1990attention,bundesen1990theory,desimone1995neural,corbetta2002control,petersen2012attention}. Ces \'etudes ont \'et\'e une source d'inspiration pour plusieurs mod\`eles d'intelligence artificielle \citep{khosla2007bio,lindsay2018biological,vaishnav2021understanding,vaishnav2022mareo}. 

Il existe trois cat\'egories de s\'electivit\'e dans un syst\`eme d'attention visuelle: par localisation spatiale \textit{(space-based)}~\citep{posner1980orienting,posner1982neural}, par appartenance \`a un objet \textit{(object-based)}~\citep{duncan1984selective,egly1994shifting,vecera1994does,kramer1997object} et par des caract\'eristiques particuli\`eres de l'entr\'ee \textit{(feature-based)}~\citep{harms1983color,driver1989movement,kramer1991perceptual,baylis1992visual,duncan1996objects}. 

L'attention est \'egalement sollicit\'ee lors de l'ex\'ecution de tâches n\'ecessitant des signaux sensoriels multiples. En pr\'esence de tâches ou de signaux sensoriels multiples, le contrôleur ex\'ecutif central aide \`a diriger l'attention. Le contrôleur ex\'ecutif central est charg\'e de coordonner l'activit\'e avec le syst\`eme cognitif pour diriger l'attention, prendre des d\'ecisions et maintenir les objectifs de la tâche. Le contexte et l'historique sont consid\'er\'es comme utiles \`a l'ex\'ecution optimale des tâches, ce qui les rend \'etroitement li\'es \`a la m\'emoire de travail. L'attention est en outre consid\'er\'ee comme le r\'esultat du contrôleur central. Le contrôleur s\'electionne les cibles de l'attention et les transmet au syst\`eme responsable de sa mise en œuvre. Il existe une relation tripartite entre le contrôle ex\'ecutif, la m\'emoire de travail et l'attention, de telle sorte que le centre d'attention est s\'electionn\'e par le contrôleur ex\'ecutif en fonction du contenu de la m\'emoire de travail ~\citep{soto2008automatic}. Bien que tous les objets de la m\'emoire de travail puissent influencer l'attention, le contrôleur ex\'ecutif aide \`a d\'ecider lequel doit affecter le plus~\citep{olivers2011difference}. Ces \'etudes cognitives vastes et approfondies li\'ees \`a l'attention ont inspir\'e le domaine de l'IA et ont contribu\'e \`a stimuler ses performances.

L'attention a profond\'ement marqué le domaine de la vision par ordinateur et du traitement naturel du langage (NLP), qui a connu un essor des architectures bas\'ees sur la \textit{self-attention} atteignant des performances de pointe sur de nombreux benchmarks. En outre, l'attention est un processus cognitif qui permet de se concentrer sur un stimulus pertinent. Cette caract\'eristique joue un rôle essentiel dans l'enrichissement de la capacit\'e de raisonnement des humains. Pour mieux comprendre le m\'ecanisme d'auto-attention, cette th\`ese a \'etudi\'e son rôle dans les architectures cognitives et de vision par ordinateur sous l'angle du raisonnement visuel.    

\color{black}
Le raisonnement visuel est le processus d'analyse des informations visuelles à disposition afin de r\'esoudre une tâche. Il est consid\'er\'e comme une partie importante de l'intelligence fluide, qui implique de penser et de raisonner ind\'ependamment de l'apprentissage, de l'\'education et de l'exp\'erience. Cette capacit\'e a \'et\'e d\'emontr\'ee non seulement chez les primates \citep{gentner2021learning} mais aussi chez les abeilles \citep{giurfa2001concepts} et chez les canetons nouveau-n\'es \citep{martinho2016ducklings}. Au contraire, des \'etudes ant\'erieures \citep{Puebla2021.04.06.438551,kim2018not,ricci37same,messina2021solving} (y compris nos propres travaux) ont montr\'e que les r\'eseaux neuronaux modernes peinent \`a r\'esoudre des tâches de raisonnement visuel simples lorsqu'ils sont test\'es sur un rep\`ere populaire appel\'e test de raisonnement visuel synth\'etique (SVRT) par \citet{fleuret2011comparing}, autrement simple pour les humains. Nous avons constat\'e une tendance similaire lorsque nous avons test\'e des architectures de raisonnement populaires de types \textit{Relational Network}~\citep{santoro2017simple}, \textit{Transformer}~\citep{vaswani2017attention} ou encore \textit{ESBN}~\citep{Webb2021EmergentST} sur la tâche de raisonnement abstrait (ART) où le stimulus contient un simple caract\`ere Unicode. Par cons\'equent, le raisonnement visuel est devenu un sujet de recherche de plus en plus populaire ces derni\`eres ann\'ees , en particulier avec l'\'emergence de nombreux tests d'intelligence fluide pour des algorithmes d'IA;  notamment les tests de raisonnement visuel compositionnel (CVR) \citep{zerroug2022benchmark}, Ravens (RPM) \citep{barrett2018measuring,zhang2019raven} et les matrices visuelles progressives (V-PROM) \citep{barrett2018measuring,teney2020v} ainsi qud les Corpus de Raisonnement Abstrait (ARC) \citep{chollet2019measure}. 

L'objectif de la premi\`ere \'etude de cette th\`ese \' a été de mettre en lumi\`ere les m\'ecanismes computationnels qui sous-tendent le raisonnement visuel \`a l'aide du \textit{Synthetic Visual Reasoning Test} (SVRT)~\citep{fleuret2011comparing}. Ce défi comprend vingt-trois problèmes de classification binaire, comprenant une vari\'et\'e de tâches de raisonnement identique-diff\'erent et spatial.  Dans notre exp\'erience, nous avons syst\'ematiquement \'evalu\'e la capacit\'e d'une batterie de $N=15$ r\'eseaux neuronaux convolutifs profonds (\textit{ResNets}) -- variant en profondeur et entraîn\'es en utilisant diff\'erentes tailles d'ensembles d'entraînement -- \`a r\'esoudre chacun des probl\`emes SVRT. Nous avons trouv\'e une gamme de précision sur l'ensemble des vingt-trois tâches. Des r\'eseaux peu profonds ont facilement appris certaines tâches, et des ensembles d'entraînement relativement petits et certaines tâches ont \'et\'e difficilement r\'esolues avec des r\'eseaux beaucoup plus profonds et des ordres de grandeur plus \'elev\'es d'exemples d'entraînement. 


Sous l'hypoth\`ese que la complexit\'e de calcul des tâches individuelles peut \^etre correctement caract\'eris\'ee par le motifs de pr\'ecision des tests de ces $N=15$ r\'eseaux neuronaux, nous avons form\'e des vecteurs de pr\'ecision \`a N dimensions pour chaque tâche et ex\'ecut\'e un algorithme de regroupement hi\'erarchique. L'analyse en r\'esultant \textcolor{black}{sugg\`ere} une taxonomie des tâches de raisonnement visuel: au-del\`a de deux clusters primaires correspondant aux jugements \textit{identique-diff\'erent} (SD) vs \textit{relation spatiale} (SR), nous avons \'egalement identifi\'e une organisation plus fine avec des sous-groupes refl\'etant la nature et le nombre de relations utilis\'ees pour composer les r\`egles d\'efinissant la tâche. Nos r\'esultats sont coh\'erents avec les travaux ant\'erieurs de \citet{kim2018not}, qui a \'et\'e le premier \`a identifier une dichotomie entre les tâches SD et SR. Nos r\'esultats prolongent \'egalement les travaux ant\'erieurs de \citep{fleuret2011comparing,kim2018not,yihe2019program} en \textcolor{black}{proposant} une taxonomie plus fine des tâches de raisonnement visuel. La pr\'ecision des r\'eseaux neuronaux se refl\`ete dans le nombre de relations utilis\'ees pour d\'efinir les r\`egles de base, ce qui est attendu, mais m\'erite un examen plus approfondi. 


\color{black}
\citet{kim2018not} ont pr\'ec\'edemment sugg\'er\'e que les tâches de SD "tendent" les r\'eseaux neuronaux convolutifs. C'est-\`a-dire que, bien qu'il soit possible de trouver une architecture de r\'eseau d'une profondeur suffisante (ou un nombre d'unit\'es suffisant) qui de plus, peut r\'esoudre une version de la tâche avec une certaine configuration de stimuli (par ex, en forçant tous les stimuli \`a \^etre contenus dans une fen\^etre $\Delta H \times \Delta W$), il est aussi relativement facile de rendre la m\^eme tâche non apprenable par cette m\^eme architecture de r\'eseau \textcolor{black}{passé} un certain nombre de stimuli de configuration (par exemple, en augmentant la taille de la fen\^etre qui contient tous les stimuli). Tout se passe comme si ces r\'eseaux convolutifs \'etaient capables d'apprendre la tâche quand le nombre de stimuli reste inf\'erieur \`a leurs capacit\'es de m\'emoire, et \'echouent au-del\`a. La question de savoir si des alternatives non convolutionnelles aux réseaux de neurones convolutionnels (CNNs) test\'es ici, comme les d\'esormais populaires r\'eseaux \textit{Transformer}  \citep{Dosovitskiy2020-iq,touvron2021training,tolstikhin2021mlp}, pourraient apprendre \`a r\'esoudre plus efficacement les tâches SVRT les plus difficiles reste encore une question ouverte. Comme exp\'erience initiale, nous avons tent\'e d'entraîner et de tester un \textit{Transformer} visuel (ViT) \footnote[1]{\url{https://github.com/facebookresearch/dino}}  \citep{Dosovitskiy2020-iq} contraint d'avoir un nombre de param\`etres (21M) égal au nombre de paramètres du modèle ResNet-50 utilis\'e ici. Pour la plupart des tâches SVRT difficiles, nous n'avons pas réussi \`a obtenir avec ces transformeurs visuels des r\'esultats au-delà de ceux obtenus par les réseaux de type ResNet et ce m\^eme en utilisant 100 000 \'echantillons (le même résultat a été montré \'egalement dans \citet{messina2021recurrent}). Il convient de noter qu'un jeu de donn\'ees de 100 000 \'echantillons reste relativement petit par rapport aux standards actuels en vision, puisque ViT a été entraîné \`a partir de z\'ero.


On peut d\'emontrer que sous certaines contraintes architecturales, les perceptrons multicouches et les r\'eseaux neuronaux convolutifs, y compris les réseaux ResNets ainsi que d'autres architectures, sont des approximateurs universels. En d'autres termes, ces réseaux peuvent apprendre des correspondances arbitraires entre des images et des labels de classe. En fonction de la complexit\'e de la correspondance, on peut avoir besoin d'un nombre croissant de neurones cach\'es afin de permettre une expressivit\'e suffisante du r\'eseau ; mais si l'on dispose de suffisamment de profondeur et d'une quantit\'e suffisante d'exemples d'entraînement, les CNNs profonds peuvent apprendre des tâches de raisonnement visuel arbitraires. Bien que nous ne puissions pas spécifiquement faire d'affirmation forte pour les architectures ResNet utilis\'ees dans cette \'etude (car la preuve d'approximation universelle a été faite dans le cadre d'une seule couche sans max pooling ni batch normalization \citep{lin2018resnet}), nous avons constat\'e empiriquement que toutes les tâches SVRT pouvaient effectivement \^etre apprises pour des r\'eseaux de profondeur suffisante et \`a condition d'avoir une quantit\'e suffisante d'exemples d'entraînement.
Cependant, les CNNs profonds sont g\'en\'eralement d\'epourvus de nombreuses fonctions cognitives humaines, telles que l'attention et la m\'emoire de travail. De telles fonctions sont susceptibles de fournir un avantage considérable \`a un apprenant pour r\'esoudre certaines de ces tâches \citep{10.7551/mitpress/1187.001.0001}. Au lieu de ces fonctions cognitives, les CNNs s'appuyeraient sur leur aspect d'approximateurs universels, conduisant \`a une solution de type "force brute" qui serait moins g\'en\'erale. 
Dans ces conditions, une question ouverte est de savoir si les tâches de SVRT d\'eriv\'es de nos analyses bas\'ees sur les CNNs serait effectivement des tâches valables pour une \'etude sur des sujets humains. De plus, la pr\'ediction faite par \citet{kim2018not} \`a l'aide des CNNs et selon laquelle les tâches SD sont plus difficiles à résoudre que les tâches SR et donc qu'elles peuvent exiger des calculs suppl\'ementaires (par le biais de processus de feedback) tels que l'attention et/ou la m\'emoire de travail a \'et\'e valid\'ee exp\'erimentalement avec succ\`es par \citet{AlamiaENEURO.0267-20.2020} en utilisant des données EEG. 


\color{black}
Une preuve suppl\'ementaire des avantages des m\'ecanismes de r\'etroaction pour le raisonnement visuel a \'et\'e fournie par \citet{linsley2018learning} qui a montr\'e que les tâches de traçage de contours qui peuvent \^etre r\'esolues efficacement avec une seule couche d'un CNN r\'ecurrent peuvent n\'ecessiter plusieurs ordres de grandeur d'\'etapes de traitement suppl\'ementaires dans un CNN non r\'ecurrent pour r\'esoudre la m\^eme tâche. Cela se traduit en fin de compte par une efficacit\'e d'\'echantillonnage beaucoup plus grande pour les CNN r\'ecurrents sur les tâches de segmentation d'images naturelles \citep{linsley2020recurrent}. \textcolor{black}{La tâche \'etroitement li\'ee de l'inutilit\'e a \'egalement \'et\'e \'etudi\'ee par \citet{villalobos2021neural} qui a d\'emontr\'e l'incapacit\'e des CNN \`a apprendre une solution g\'en\'erale pour cette classe de probl\`emes.}
Les approximateurs universels avec des biais inductifs minimaux tels que les perceptrons multicouches, les CNN et d'autres architectures feedforward ou \textcolor{black}{non}-attentives peuvent apprendre \`a r\'esoudre des tâches de raisonnement visuel, mais ils peuvent avoir besoin d'un tr\`es grand nombre d'exemples d'entraînement pour s'adapter correctement. Par cons\'equent, au-del\`a de la simple mesure de la pr\'ecision de r\'eseaux tr\`es profonds dans des r\'egimes de donn\'ees \'elev\'es (comme lorsque des millions d'exemples d'entraînement sont disponibles), l'\'evaluation syst\'ematique des performances de r\'eseaux neuronaux de diff\'erentes profondeurs et pour diff\'erents r\'egimes d'entraînement peut fournir des informations essentielles sur la complexit\'e de diff\'erentes tasks de raisonnement visuel.

Plus tôt, \citet{kim2018not} a \'emis l'hypoth\`ese que cette tension des r\'eseaux convolutifs est due \`a leur manque de m\'ecanismes d'attention permettant de lier explicitement les r\'egions de l'image aux objets mentaux. Une remarque similaire a \'et\'e faite par \citet{greff2020binding} dans le contexte de l'incapacit\'e des r\'eseaux neuronaux contemporains \`a d\'ecouper les informations sensorielles en morceaux distincts qui peuvent ensuite \^etre analys\'es et compar\'es individuellement (voir \'egalement \citet{10.1007/978-3-540-75555-5_15} pour un point similaire). Il est int\'eressant de noter que cette pr\'ediction a r\'ecemment \'et\'e test\'ee \`a l'aide de l'EEG humain par \citet{AlamiaENEURO.0267-20.2020} qui a montr\'e qu'en effet l'activit\'e c\'er\'ebrale enregistr\'ee pendant les tâches SD est compatible avec des demandes d'attention et de m\'emoire de travail plus importantes que les tasks SR. En m\^eme temps, le fait que les CNN puissent apprendre des tasks SR plus efficacement que des tasks SD ne signifie pas n\'ecessairement que les participants humains peuvent r\'esoudre ces tâches sans attention. En effet, \citep{logan1994spatial} a montr\'e que les tâches SR telles que le jugement de l'int\'erieur n\'ecessitent de l'attention dans certaines circonstances.

Pour \'evaluer le rôle de l'attention dans le raisonnement visuel, nous avons utilis\'e des modules \textit{Transformer} pour doter les CNNs profonds d'une attention spatiale et d'une attention bas\'ee sur les caract\'eristiques. Les am\'eliorations relatives obtenues par les CNNs avec les deux formes d'attention varient selon les tasks. De nombreuses tasks refl\'etaient une plus grande am\'elioration de l'attention spatiale, et un plus petit nombre b\'en\'eficiait de l'attention bas\'ee sur les caract\'eristiques. De plus, nous avons constat\'e que les mod\`eles d'am\'eliorations relatives expliquaient une grande partie de la variance dans l'espace des tasks SVRT d\'eriv\'ees dans l'exp\'erience 1. Dans l'ensemble, nous avons constat\'e que l'exigence d'une attention spatiale et d'une attention bas\'ee sur les caract\'eristiques explique bien la taxonomie des tasks de raisonnement visuel identifi\'ees dans l'exp\'erience 1.  Notre analyse informatique a \'egalement conduit \`a des pr\'edictions testables pour les exp\'eriences humaines en sugg\'erant des tasks qui b\'en\'eficient soit de l'attention spatiale (tâche \textit{22}) soit de l'attention bas\'ee sur les caract\'eristiques (tâche \textit{21}), des tâches qui b\'en\'eficient des deux formes d'attention (tâche \textit{19}) et des tâches qui ne b\'en\'eficient pas de l'attention (tâche\textit{2}).


\color{black}
Enfin, notre \'etude s'est concentr\'ee sur les avantages computationnels de l'attention spatiale et de l'attention bas\'ee sur les caract\'eristiques pour le raisonnement visuel. Les travaux futurs devraient consid\'erer le rôle d'autres formes d'attention, y compris l'attention bas\'ee sur l'objet \citep{egly1994covert} pour le raisonnement visuel.

Dans notre deuxi\`eme exp\'erience, nous avons \'etudi\'e la capacit\'e d'apprentissage des caract\'eristiques SVRT par rapport aux r\`egles. Pour ce faire, nous avons pr\'eform\'e les r\'eseaux neuronaux sur des tasks auxiliaires afin d'apprendre les caract\'eristiques SVRT avant de les former \`a apprendre les r\`egles abstraites associ\'ees aux probl\`emes SVRT individuels. Nos m\'ethodes de pr\'e-entraînement ont conduit \`a des r\'eseaux qui apprennent \`a r\'esoudre les probl\`emes SVRT mieux que les r\'eseaux form\'es \`a partir de z\'ero, ainsi que les r\'eseaux qui ont \'et\'e pr\'e-entraîn\'es pour effectuer la cat\'egorisation d'images sur le jeu de donn\'ees ImageNet. Nous avons \'egalement constat\'e que ces processus d'attention semblent contribuer davantage \`a l'apprentissage de r\`egles qu'\`a l'apprentissage de caract\'eristiques. Pour le sous-cluster $SR_1$, nous constatons que ce type de pr\'e-entraînement est avantageux dans les r\'egimes d'entraînement inf\'erieurs, mais que les avantages disparaissent rapidement dans les r\'egimes d'entraînement sup\'erieurs. En revanche, ce pr\'e-entraînement ne permet pas d'apprendre les tasks du sous-cluster $SD_1$, m\^eme avec 15 000 \'echantillons, ce qui sugg\`ere que le principal d\'efi de ces tasks n'est pas de d\'ecouvrir de bonnes repr\'esentations visuelles, mais plutôt de d\'ecouvrir la r\`egle. Cela sugg\`ere le besoin de m\'ecanismes suppl\'ementaires au-del\`a de ceux mis en œuvre dans les ResNets. Ceci est \'egalement coh\'erent avec les am\'eliorations observ\'ees pour ces tasks avec l'ajout de m\'ecanismes d'attention. 

En r\'esum\'e, notre \'etude a compar\'e les exigences informatiques de diff\'erentes tasks de raisonnement visuel. Bien que nous nous soyons concentr\'es sur la compr\'ehension des avantages computationnels des m\'ecanismes d'attention et d'apprentissage des caract\'eristiques, il est clair que des m\'ecanismes suppl\'ementaires seront n\'ecessaires pour r\'esoudre pleinement toutes les tasks de SVRT. Ces m\'ecanismes sont susceptibles d'inclure la m\'emoire de travail, dont on sait qu'elle joue un rôle dans les tasks de SVRT \citep{AlamiaENEURO.0267-20.2020}. Dans l'ensemble, ce travail illustre les avantages potentiels de l'incorporation de m\'ecanismes semblables \`a ceux du cerveau dans les r\'eseaux neuronaux modernes et fournit une voie \`a suivre pour atteindre un raisonnement visuel de niveau humain. 

Dans la derni\`ere partie de la th\`ese, nous avons propos\'e une nouvelle architecture, le Guided Attention Model for (visual) Reasoning (\textit{GAMR}). Nous avons int\'egr\'e les deux capacit\'es cognitives utilis\'ees par les humains - l'attention et la m\'emoire - dans la r\'esolution de tasks de raisonnement. Les th\'eories cognitives modernes de la vision active postulent que le syst\`eme visuel explore l'environnement de mani\`ere dynamique par le biais de s\'equences de changements d'attention pour s\'electionner et acheminer vers la m\'emoire les informations pertinentes pour la task \citep{ullman1984massachusetts,ullman1987visual}. Des exp\'eriences de psychophysique \citep{hayhoe2000vision} sur l'attention visuelle manifeste ont montr\'e que les sch\'emas de mouvements oculaires sont dirig\'es selon des routines d\'ependantes de la task. Le GAMR s'inspire de la litt\'erature des sciences cognitives sur la vision active, où l'attention manifeste est dirig\'ee dans le syst\`eme visuel pour recueillir des informations pertinentes pour la task. Selon la th\'eorie de la vision active, le monde visuel est explor\'e \`a l'aide de mouvements oculaires rapides guid\'es par des d\'eplacements de l'attention visuelle. 


\color{black}
Nous avons conçu un contrôleur semblable aux m\'ecanismes impliqu\'es dans le cadre de la vision active pour diriger le projecteur d'attention et envoyer les repr\'esentations pertinentes pour la task dans le bloc de m\'emoire utilis\'e ensuite pour le raisonnement. Dans GAMR, le contrôleur est impl\'ement\'e avec la couche d'encodeur bas\'ee sur les transformateurs. Contrairement \`a la m\'ethode existante où l'auto-attention est suivie d'une addition et d'une normalisation de la couche et où une couche lin\'eaire est ajout\'ee \`a la fin de cette op\'eration, nous supprimons la couche d'auto-attention et la remplaçons par l'attention bas\'ee sur les caract\'eristiques obtenue avec le module contrôleur LSTM. Cela aide le contrôleur \`a d\'eplacer le point d'attention. 
 
Notre proposition de module d'attention guid\'ee pour le raisonnement (visuel) (GAMR) apprend \`a d\'eplacer l'attention dynamiquement, en fonction de la task, sur la base de requ\^etes g\'en\'er\'ees en interne par un contrôleur ex\'ecutif LSTM. Grâce \`a des exp\'eriences approfondies sur les deux principaux d\'efis de raisonnement visuel, le Synthetic Visual Reasoning Test (SVRT)~\citep{fleuret2011comparing} et le Abstract Reasoning Task (ART)~\citep{Webb2021EmergentST}, nous d\'emontrons que notre architecture neuronale est capable d'apprendre des compositions complexes de r\`egles relationnelles d'une mani\`ere efficace en termes de donn\'ees et qu'elle est plus performante que les autres architectures neuronales de pointe pour le raisonnement visuel. En utilisant des m\'ethodes d'explicabilit\'e, nous caract\'erisons davantage les strat\'egies visuelles utilis\'ees par le mod\`ele afin de r\'esoudre des tasks de raisonnement repr\'esentatives. Nous d\'emontrons que notre mod\`ele est compositionnel, c'est-\`a-dire qu'il est capable de se g\'en\'eraliser efficacement \`a de nouvelles tasks et d'apprendre de nouvelles routines visuelles en recomposant des op\'erations \'el\'ementaires apprises pr\'ec\'edemment. 

Nous nous sommes inspir\'es pour la banque de m\'emoire de \citet{Webb2021EmergentST}, où des m\'ecanismes de liaison et d'indirection de variables ont \'et\'e introduits dans l'architecture pour le raisonnement visuel \`a l'aide d'une m\'emoire externe. La liaison de variables est la capacit\'e de lier deux repr\'esentations, et l'indirection est le m\'ecanisme impliqu\'e dans la r\'ecup\'eration d'une repr\'esentation pour se r\'ef\'erer \`a l'autre. Ces auteurs introduisent \'egalement la normalisation temporelle du contexte (TCN) \citep{webb2020learning}, qui s'av\`ere b\'en\'efique pour la g\'en\'eralisation hors distribution pour les tasks de raisonnement relationnel. Cependant, le mod\`ele pr\'esente des limites importantes: Il suppose une repr\'esentation d'image centr\'ee sur l'objet, où les objets sont pr\'esent\'es individuellement dans une s\'equence. Nous ne pouvons pas \'evaluer une telle architecture dans le cadre du d\'efi SVRT, car les images de chaque task contiennent plusieurs objets qui n\'ecessitent une individualisation. Il existe \'egalement certaines relations, comme "toucher", que cette individuation ne peut repr\'esenter (\textcolor{black}{ou toute architecture centr\'ee sur l'objet}). ESBN n'a pas non plus de m\'ecanisme attentionnel et fonctionne mieux dans un sc\'enario où une attention forte au niveau du pr\'e-traitement permet de simplifier les tasks. Nous avons test\'e ce comportement de correspondance des mod\`eles de l'architecture en l'entraînant en pr\'esence d'un bruit gaussien. Il a conduit \`a une performance de niveau al\'eatoire. Ici, nous nous appuyons sur ce travail et d\'ecrivons un mod\`ele entraînable de bout en bout qui apprend \`a individualiser les sc\`enes pertinentes pour la task et \`a stocker leurs repr\'esentations en m\'emoire pour permettre de juger des relations complexes entre ces objets. Enfin, notre m\'ecanisme relationnel est inspir\'e du travail de \citet{santoro2017simple} qui a introduit un mod\`ele pr\^et \`a l'emploi pour calculer les relations entre les repr\'esentations de type objet dans un r\'eseau. 


\color{black}

L'une des limites de l'approche actuelle est le nombre fixe de pas de temps. Je pense qu'une future continuation de ce travail pourrait \^etre d'incorporer un m\'ecanisme pour adapter le nombre de pas de temps en fonction de la complexit\'e de la task. Pour l'instant, nous avons fix\'e le nombre de pas de temps \`a quatre pour toutes les tasks ; cependant, une task simple pourrait n\'ecessiter moins de pas de temps pour arriver \`a une d\'ecision avec une grande confiance. Pour que le mod\`ele s'adapte \`a la situation, une possibilit\'e pourrait \^etre de l'entraîner avec une variable de confiance comme crit\`ere d'arr\^et. 

Bien que nous ayons limit\'e notre analyse \`a des ensembles de donn\'ees synth\'etiques de raisonnement visuel, il serait possible \`a l'avenir de tester les mod\`eles sur un ensemble de donn\'ees du monde r\'eel comme V-PROM. Il s'agit d'images organis\'ees dans le style de raisonnement des Ravens, avec des images de contexte et des images de choix parmi lesquelles la bonne r\'eponse est s\'electionn\'ee. Une autre direction possible est de penser \`a une architecture qui prend en compte deux caract\'eristiques importantes: l'utilisation efficace des donn\'ees et l'utilisation efficace des ressources informatiques. Une façon de concevoir cette architecture est d'incorporer un m\'ecanisme de lecture et d'\'ecriture similaire \`a une machine de Turing neuronale. Ces deux m\'ecanismes aideront le r\'eseau \`a lire les relations d\'ej\`a stock\'ees en m\'emoire et \`a les \'ecrire dans la m\'emoire si elles sont nouvelles. Nous nous attendons \`a ce qu'une telle architecture cognitive d\'emontre une capacit\'e de raisonnement d'ordre sup\'erieur, d'apprentissage continu, de composition et de m\'eta-apprentissage.

Nous avons \'egalement \'evalu\'e ViT~\citep{Dosovitskiy2020-iq} -- Nous avons \'egalement \'evalu\'e ViT~citep{Dosovitskiy2020-iq} - une architecture d'auto-attention compl\`ete sur des tasks SVRT et avons constat\'e qu'elle avait du mal \`a apprendre les tasks SVRT les plus simples ; cependant, \citet{messina2021recurrent} a men\'e une \'etude similaire sur un sous-ensemble plus petit de quatre tasks SVRT entra\^in\'ees sur 28k \'echantillons et a constat\'e qu'une version r\'ecurrente de ViT - un r\'eseau attentionnel avec une colonne vert\'ebrale convolutive peut apprendre ces tasks. L'ajout de convolutions dans les premi\`eres couches de ViT permet d'obtenir une meilleure pr\'ecision et am\'eliore la sensibilit\'e aux param\`etres d'optimisation \citep{xiao2021early}. Cette observation nous a motiv\'e \`a proposer \textit{Conviformer} \citep{vaishnav2022conviformers} pour un autre projet collaboratif sur la classification des feuilles-fossiles. Nous proposons un r\'eseau pour incorporer un r\'eseau convolutif comme frontal d'un r\'eseau de transformation de la vision bas\'e sur l'auto-attention compl\`ete, am\'eliorant ainsi sa capacit\'e \`a traiter des images \`a plus haute r\'esolution. Bien que les images de plus grande taille rev\^etent une grande importance dans les applications de vision par ordinateur telles que la d\'etection d'objets, la segmentation et la classification \`a grain fin, elles ne peuvent pas \^etre utilis\'ees avec les transformateurs de vision en raison de la demande de m\'emoire de calcul associ\'ee. \textit{Conviformer} am\'eliore la performance des transformateurs de vision en incorporant des caract\'eristiques locales et en infusant des prieurs convolutifs dans une architecture de transformateur. Nous aimerions voir comment les transformateurs de vision induits par convolution se comportent dans les tasks SVRT. 

L'apprentissage de concepts est une autre direction de recherche passionnante. L'une des principales caract\`eristiques de l'intelligence humaine est la capacit\'e d'apprendre rapidement de nouveaux concepts et de les utiliser pour g\'en\'eraliser \`a un nouveau sc\'enario. Un \textit{concept} peut \^etre une id\'ee repr\'esentant une classe d'\'ev\'enements (par exemple, la marche), des objets (par exemple, les chats), ou leurs propri\'et\'es (par exemple, la couleur bleue). Pour tester la capacit\'e d'apprentissage de concepts des r\'eseaux neuronaux en quelques coups, nous avons r\'ecemment introduit un nouveau jeu de donn\'ees de raisonnement visuel, Compositional Visual Reasoning (\textit{CVR})~\citep{zerroug2022benchmark}. Ce jeu de donn\'ees est bas\'e sur le principe du raisonnement odd-one-out. Dans cette forme de task de raisonnement, trois \'echantillons sur quatre suivent un concept similaire (r\`egle) dans leur formulation, tandis que le quatri\`eme ne le fait pas. Chaque \'echantillon contient des formes similaires aux formes utilis\'ees dans le d\'efi SVRT. Cela permet d'\'etendre la vari\'et\'e des relations utilis\'ees dans la formulation par rapport aux jeux de donn\'ees pr\'ec\'edemment d\'efinis comme SVRT ou RPM. Nous avons \'egalement inclus la compositionnalit\'e pr\'ealable dans le jeu de donn\'ees, où certaines relations \'el\'ementaires sont utilis\'ees pour composer les diff\'erentes tasks. La motivation est de pousser la communaut\'e \`a construire un r\'eseau compositionnel et efficace en termes d'\'echantillons.

Dans cette th\`ese, nous avons fait l'une des toutes premi\`eres tentatives d'explorer l'auto-attention du point de vue du raisonnement visuel. L'attention joue un rôle crucial dans la d\'emonstration des capacit\'es de raisonnement visuel, et on s'attend \`a ce qu'un meilleur mod\`ele attentionnel soit meilleur pour le raisonnement. Nous avons montr\'e comment les op\'erations d'auto-attention pouvaient \^etre utilis\'ees comme mod\`ele computationnel d'un syst\`eme d'attention visuelle repr\'esentant l'attention spatiale et bas\'ee sur les caract\'eristiques, ainsi que comme mod\`ele de vision active. Bien que nous ayons constat\'e que l'auto-attention est aussi efficace dans la r\'esolution de tasks de raisonnement que dans d'autres d\'efis li\'es \`a la vision, il est n\'ecessaire d'effectuer des analyses suppl\'ementaires pour comprendre les m\'ecanismes fondamentaux d'un mod\`ele complet d'auto-attention qui limitent son apprentissage efficace par \'echantillonnage pour les tasks de raisonnement. Dans l'ensemble, ce travail d\'emontre les avantages potentiels de l'ajout de m\'ecanismes d'auto-attention dans l'architecture cognitive et de vision par ordinateur pour r\'esoudre les tasks de raisonnement visuel. 